%% sections/02-related-work.tex

\section{Related Work}
\label{sec:related}

\subsection{Resource Management in Games}

Resource management is a core category of game mechanics, studied extensively in
strategy games~\cite{juul2005half}, roleplaying games~\cite{dungeons2014players},
and survival games~\cite{fullerton2008}. The standard model is that resource pressure
creates meaningful choices by forcing players to allocate finite resources across
competing demands~\cite{salen2004rules}. This produces difficulty and engagement with
the game's mechanical system.

\citet{schell2008art} describes resource management as a lens for creating tension:
``if the player never runs out of anything, there's no scarcity, and without scarcity,
there's no tension.'' This is the tactical model: tension comes from scarcity of
action-enabling resources.

We add an alternative: tension can also come from scarcity of relationship-enabling
resources. When the resource represents people rather than power, its depletion
creates narrative tension (who will we lose?) rather than only tactical tension
(will we survive?). Narrative tension activates character-driven decision-making;
tactical tension activates combat-optimization thinking.

\subsection{Attrition in Long-Form Games}

Long-form game design faces the \textit{attrition problem}: in campaigns with many
sessions, early resource expenditures are often forgotten by players by the time they
matter~\cite{perkins2017dungeon}. The Marathon mechanic addresses this with explicit
tracking (Supply Die die-size; Cohesion Die die-size) that makes cumulative attrition
visible at each session.

\citet{alexandrian2012} describes the ``three-clue rule'' for mystery scenarios as
a reliability mechanism: if one delivery of a clue fails, two backups ensure delivery.
The Supply Die's explicit degradation track is an analogous reliability mechanism for
attrition: even if a player doesn't feel the HP carry-over between sessions, the Supply
Die's d12$\to$d10$\to$d8$\to$d6$\to$d4$\to$exhausted progression makes cumulative
attrition undeniable.

\subsection{Player Motivation and Resource Pressure}

\citet{yee2006motivations} documents that MMORPG players with high Achievement motivation
engage more intensely with resource-pressure mechanics. \citet{nacke2011brainhex}
finds that Conqueror and Seeker player types respond to scarcity differently. Our
data cannot distinguish player motivation (all sessions are AI-simulated) but can
distinguish resource-mechanic design type (Supply Die vs.\ Cohesion Die) and
provide Character Agency scores as proxies for character-driven engagement.

\subsection{Gap in the Literature}

No prior work, to our knowledge, compares resource mechanics on their
\textit{narrative amplification} effect (Character Agency scores) rather than their
tactical-difficulty effect. The distinction matters for narrative game design:
a designer who wants players to make harder tactical choices should use Supply Die
mechanics; a designer who wants players to make deeper character-driven choices
should use Cohesion Die mechanics. Both produce resource pressure; only the Cohesion
Die strongly amplifies the character-driven dimension in sessions 5--7.
